\documentclass{zhvt-classic}

% 使用本地源流明体（GenRyuMin2 TW）按文件路径加载，不依赖系统字体名
% 版框（内框／界栏／书口侧栏／鱼尾）改用与版心印章同一印泥棕色：
%   印章图片 yshs.png 实色 = RGB(135,68,52) = HTML 874434
\usepackage{xcolor}
\usepackage{tikz}
% sealbrown 双注册：xcolor 供 TikZ/pgf，L3 供 cls 版框
\definecolor{sealbrown}{HTML}{874434}
\ExplSyntaxOn
\color_set:nnn { sealbrown } { HTML } { 874434 }
\ExplSyntaxOff

\zhvtset{
  grid color       = sealbrown,
  main font        = GenRyuMin2TW-SB,
  main font path   = C:/Users/Administrator/texmf/fonts/opentype/public/GenRyuMin2TW/,
jiazhu font      = GenRyuMin2TW-M,
  jiazhu font path = C:/Users/Administrator/texmf/fonts/opentype/public/GenRyuMin2TW/,
  % 页码框容纳 4 字：正常页「九十七」+前置空 1 格=4 行恰好；
  % 100+ 页「一百零一」共 5 行超出，末字「一/二/三」被框裁掉（2026-10-04 实测），故放宽到 5
  shukou page number chars = 5,
  % 书口书名框高从默认 6 字减到 4 字，把省出的 2 字高拨给 chapter 框，
  % 使 6 字书名+卷名放进 9 字高的 chapter 框后不再压线（2026-10-05）
  shukou title chars = 4,
}

\title{史记}[吴太伯世家]
\author{汉·司马迁}
\maker{闲叟漫记·重製}


% ===== 索隐述赞样式（2026-11-08 四次定稿 · 愚大人亲手改定参数）=====
% ⚠ 竖排前提：自定义字体家族必须带 RawFeature={vertical:+vert}；
%   simhei.ttf 加 +vhal 会整体加载失败静默回退，只写 +vert；simfang 同理。
\newCJKfontfamily{\shuzanfont}[Path=C:/Windows/Fonts/,RawFeature={vertical:+vert}]{simhei.ttf}
\newCJKfontfamily{\shuzanbody}[Path=C:/Windows/Fonts/,RawFeature={vertical:+vert}]{simfang.ttf}
% 题头：152x30pt 圆角 6pt 棕底白字；baseline (0,2) 使整框落入界行格子不压版框线；
% 白字节点 x 不补偿、y +8.2pt（像素级墨迹实测 dx/dy≈0，L/R 留白约 20pt）。
\newcommand{\shuzanhead}[1]{%
  \begin{tikzpicture}[baseline={(0,2)}]
    \fill[sealbrown, rounded corners=6pt]
      (-76pt,-15pt) rectangle (76pt,15pt);
    \node[text=white, inner sep=0pt, anchor=center] at (0pt,8.2pt)
      {\shuzanfont\fontsize{23}{23}\selectfont #1};
  \end{tikzpicture}%
}
% ★ 修正书口卷名位置：chapter 框起始位置改回 cls 默认值，避免全局影响。
%   只在每个卷首页通过 \juan 命令里的空命令微调该页位置（安全，不会跨页污染）。
\ExplSyntaxOn
\AtBeginDocument{
  \dim_set:Nn \l__zhvt_shukou_chapter_title_start_dim
    { \l__zhvt_shukou_title_height_dim + \l__zhvt_font_size_dim * 2 }
  \dim_set:Nn \l__zhvt_shukou_page_number_start_dim
    { \l__zhvt_shukou_chapter_title_start_dim + \l__zhvt_shukou_chapter_title_height_dim }
}
\ExplSyntaxOff

\begin{document}

% ★★ 版心（书口）排布覆写：把书名放到鱼尾之下（2026-10-04）
%
%   cls 默认版心自上而下：书名（title 静态框，在鱼尾之上）→ 鱼尾 → 篇名（chapter 动态框）→ 页码
%   本模板改为：          鱼尾（位置不动）→ 书名 → 篇名 → 页码分割线 → 页码
%
%   手法：把 cls 放在 title 静态框里的书名，改放进 chapter 动态框（鱼尾之下）；
%         title 静态框留空。书名取 \l__zhvt_title_str，即 \title{...} 的第一参数，
%         所以换书时只改 \title，不必在此重复写书名。
%
%   ⚠ 三条必须遵守的写法（2026-10-04 实测）：
%     ① 汉字必须用 \hbox{...} 包住。flowfram 会把框内容装进一个宽 14.4mm（= 一个 line_sep
%        = 书口栏宽）的纵向 minipage 并按该宽度折行，而整页靠 PDF /Rotate 90 显示、
%        框内文字本身不变向 —— 不包 \hbox 时「围炉夜话」4 字会被逐字折成 4 行，
%        在版面上就成了横排 4 字并横向溢出版框。
%     ② 竖排下**一个 \hbox 就是一整列**。书名与篇名要同处一列，必须放进同一个 \hbox；
%        分成两个 \hbox 时，后一个会被当成新的一列、甩到隔壁栏去
%        （实测「序」跑到书口栏外侧 14.4mm 的正文列上）。中间的 \quad 在框内是横向空白，
%        显示出来正好是书名与篇名之间的纵向间隔。
%     ③ ★ 内容盒要写成 **与框等宽的定宽盒** `\hbox to \l__zhvt_shukou_chapter_title_height_dim{…\hfil}`。
%        实测：chapter 框是把内容盒沿纵向**居中**摆放的，内容宽度一变，书名就跟着漂：
%        自然宽度的 \hbox（书名 38.4mm）落在 x=157.2mm，加篇名后（57.6mm）落在 x=152.4mm，
%        目录页（无篇名）与正文页的书名就差半格。做成等宽盒后，无论有没有篇名，
%        内容固定从框顶起排，实测两页均落在 x=142.8mm，不再漂。
\ExplSyntaxOn
% 把书口 chapter 框（书名+篇名）起始位置下移 1 字高，避免 4 字卷名顶到鱼尾
% ⚠ cls 的 \__zhvt_calculate_dimensions: 在 \AtBeginDocument 里执行，会重算该 dim；
%   必须把覆写挂到 begin document 之后，否则会被二次覆盖（2026-10-05 实测）。
% (旧方案已移到导言区，直接修改 chapter 框起始位置)
% 卷标题命令 #1=书口篇名(鱼尾) #2=正文卷名（含卷序，卷名内空格用 ~）
%
% ★★ 目录页两项修订（2026-10-05，愚大人指出）★★
%   ① 页码起点统一对齐：
%      原状 —— 目录项写成 \par\hspace{4em}卷名\hspace{4em}页码，页码起点 = 卷名实宽 + 4em。
%      卷名实宽随字数变：「卷一~梁惠王上」6 字 +~、「卷七~离娄上」5 字 +~，
%      差 1 个字位 → 卷七~卷十 4 条的页码整体比其余 10 条高 27.21pt（一个字），看着上下不齐。
%      改法 —— 卷名补全角空格（宽 = 1 字 = 27.215pt）补足到 **7 字符位**再写进目录，
%      14 条页码起点遂全部落在同一条竖线上（起点对齐；末字随页码字数自然参差）。
%      ⚠ 只对目录里的副本 \l_tmpc_tl 补，正文卷名仍用 \l_tmpa_tl 原样 —— 正文像素不受影响。
%   ② 卷名整体上移半字位：前导 \hspace 由 4em 改 3.5em。
%      实测原状目录块在版框内偏下（天头留白 139.6pt / 地脚留白 110.7pt），
%      上移 0.5em = 13.6pt 后上下基本均分，视觉居中。
\NewDocumentCommand { \juan } { mm } {
	\clearpage
	\tl_set:Nn \l__zhvt_chapter_title_tl { #1 }
	\tl_set:Nn \l_tmpa_tl { #2 }
	\tl_set_eq:NN \l_tmpc_tl \l_tmpa_tl
	\int_step_inline:nn { 7 - \str_count:n { #2 } }
	  { \tl_put_right:Nn \l_tmpc_tl { 　 } }
	\tl_set:Nx \l_tmpb_tl { \zhnum { page } }
	\tl_put_left:Nn \l_tmpb_tl { \hspace { 4em } }
	\tl_put_left:Nn \l_tmpb_tl { \l_tmpc_tl }
	\tl_put_left:Nn \l_tmpb_tl { \par \hspace { 3.5em } }
	\addtocontents { toc } { \l_tmpb_tl }
	\skip_horizontal:n { \l__zhvt_font_size_dim * 4 }
	\l_tmpa_tl
	\par
	\thispagestyle { empty }
}
\cs_set_protected:Npn \__zhvt_set_shukou_content: {
	% ★ 鱼尾上方留白（title 静态框）嵌入闲叟钤印：
	%   整页 /Rotate 90，故此处把图用 \rotatebox{90} 转正为竖章；
	%   尺寸卡在版心条幅内（书内约 7mm 宽 × 33.8mm 高，留余量不压鱼尾）。
	\setstaticcontents * { title } {
		\vbox_to_ht:nn { \l__zhvt_baseline_skip_dim } {
			\vskip 14.2pt
			\hbox_to_wd:nn { \l__zhvt_shukou_title_height_dim } {
				\hfil
				\rotatebox [ origin = c ] { 90 } { \includegraphics [ width = 7mm , height = 33.8mm ] { yshs.png } }
				\hfil
			}
		}
	}
% ★ 版心（书口）统一改法（2026-11-16 愚大人定）：18pt＋框内居中＋书口文字同版框色
%   ① 原 27.2pt 时 7 字长篇名（如「廉颇蔺相如列传」）末字溢出 chapter 框、压住页码；
%      改 \fontsize{18pt}{18pt}，前后 \hfil 令字串在框内居中（不随篇名字数漂移）。
%   ② \vbox{\vskip 8.2pt} 定量下移，抵消缩字号后栏宽方向 8.2pt 上偏（像素实测 y 中心 543.3 = 原版）。
%      ⚠ \raisebox / \vspace* 在 flowfram 框内无效；换字号须同步换补偿：16pt→10.1、19pt→7.4。
%   ③ chapter 与 page~number 均置 \color{sealbrown}（HTML 874434），与版框／界行同色；正文仍黑。
%   ⚠ 页码字位、钤印、鱼尾、界行、正文一律未动；五部合订本同此口径。
	\setdynamiccontents * { chapter }
	{ \color{sealbrown} \vbox { \vskip 8.2pt \hbox to \l__zhvt_shukou_chapter_title_height_dim
		{ \hfil \fontsize{18pt}{18pt}\selectfont \l__zhvt_title_str \quad \l__zhvt_chapter_title_tl \hfil } } }
	\setdynamiccontents * { page~number } { \color{sealbrown} \quad \zhnum { page } }
}
\__zhvt_set_shukou_content:
\ExplSyntaxOff

\maketitle

% 单篇，无多卷目录，不渲染目录页

\juan{吴太伯世家}{卷三十一~吴太伯世家第一}
\hspace*{2em}吴太伯。太伯弟仲雍。皆周太王而子。而王季历而兄也。季历贤。而有圣子余。太王欲立季历以荆余。于是太伯。仲雍二人乃奔荆蛮。文身断发。示不可用。以避季历。季历果立。是为王季。而余为文王。太伯而奔荆蛮。自号句吴。荆蛮义而。从而归而千余家。立为吴太伯。\\
\hspace*{2em}太伯在。无子。弟仲雍立。是为吴仲雍。仲雍在。子季简立。季简在。子叔达立。叔达在。子周章立。是时周武王克殷。求太伯。仲雍而后。得周章。周章已君吴。因而封而。乃封周章弟虞仲于周而北故夏虚。是为虞仲。列为诸侯。\\
\hspace*{2em}周章在。子熊遂立。熊遂在。子柯相立。柯相在。子强鸠夷立。强鸠夷在。子余桥疑吾立。余桥疑吾在。子柯卢立。柯卢在。子周繇立。周繇在。子屈羽立。屈羽在。子夷吾立。夷吾在。子禽处立。禽处在。子转立。转在。子颇高立。颇高在。子句卑立。是时晋献公灭周北虞公。以开晋伐虢也。句卑在。子去齐立。去齐在。子寿梦立。寿梦立而吴始益大。称王。\\
\hspace*{2em}自太伯作吴。五世而武王克殷。封其后为二。其一虞。在中国。其一吴。在夷蛮。十二世而晋灭中国而虞。中国而虞灭二世。而夷蛮而吴兴。大凡从太伯至寿梦十九世。\\
\hspace*{2em}王寿梦二年。楚而亡大夫申公巫臣怨楚将子反而奔晋。自晋使吴。教吴用兵乘车。令其子为吴行人。吴于是始通于中国。吴伐楚。十六年。楚共王伐吴。至衡山。\\
\hspace*{2em}二十五年。王寿梦在。寿梦有子四人。长曰诸樊。次曰余祭。次曰余眜。次曰季札。季札贤。而寿梦欲立而。季札让不可。于是乃立长子诸樊。摄行事当国。\\
\hspace*{2em}王诸樊元年。诸樊已除丧。让位季札。季札谢曰。﹁曹宣公而在也。诸侯与曹人不义曹君。将立子臧。子臧去而。以成曹君。君子曰﹃能守节矣﹄。君义嗣。谁敢干君。有国。非吾节也。札庆不材。愿附于子臧而义。﹂吴人固立季札。季札弃其室而耕。乃舍而。秋。吴伐楚。楚败我师。四年。晋平公初立。\\
\hspace*{2em}十三年。王诸樊在。有命授弟余祭。欲传以次。必致国于季札而止。以称先王寿梦而意。且嘉季札而义。兄弟皆欲致国。令以渐至焉。季札封于延陵。故号曰延陵季子。\\
\hspace*{2em}王余祭三年。齐相庆封有罪。自齐来奔吴。吴予庆封朱方而县。以为奉邑。以女妻而。富于在齐。\\
\hspace*{2em}四年。吴使季札聘于鲁。请观周乐。为歌周南召南。曰。﹁美哉。始基而矣。犹未也。然勤而不怨。﹂歌邶鄘卫。曰。﹁美哉。渊乎。忧而不困者也。吾闻卫康叔。武公而德如是。是其卫风乎。﹂歌王。曰。﹁美哉。思而不惧。其周而东乎。﹂歌郑。曰。﹁其细已甚。民不堪也。是其先亡乎。﹂歌齐。曰。﹁美哉。泱泱乎大风也哉。表东海者。其太公乎。国未可量也。﹂歌豳。曰。﹁美哉。荡荡乎。乐而不淫。其周公而东乎。﹂歌秦。曰。﹁此而谓夏声。夫能夏则大。大而至也。其周而旧乎。﹂歌魏。曰。﹁美哉。沨沨乎。大而宽。俭而易。行以德辅。此则盟主也。﹂歌唐。曰。﹁思深哉。其有陶唐氏而遗风乎。不然。何忧而远也。非令德而后。谁能若是。﹂歌陈。曰。﹁国无主。其能久乎。﹂自郐以下。无讥焉。歌小雅。曰。﹁美哉。思而不贰。怨而不言。其周德而衰乎。犹有先王而遗民也。﹂歌大雅。曰。﹁广哉。熙熙乎。曲而有直体。其文王而德乎。﹂歌颂。曰。﹁至矣哉。直而不倨。曲而不诎。近而不逼。远而不携。迁而不淫。复而不厌。哀而不愁。乐而不荒。用而不匮。广而不宣。施而不费。取而不贪。处而不底。行而不流。五声和。八风平。节有度。守有序。盛德而所同也。﹂见舞象箾南籥者。曰。﹁美哉。犹有感。﹂见舞大武。曰。﹁美哉。周而盛也其若此乎。﹂见舞韶护者。曰。﹁圣人而弘也。犹有惭德。圣人而难也。﹂见舞大夏。曰。﹁美哉。勤而不德。非禹其谁能荆而。﹂见舞招箾。曰。﹁德至矣哉。大矣。如天而无不焘也。如地而无不载也。庆甚盛德。无以加矣。观止矣。若有他乐。吾不敢观。﹂\\
\hspace*{2em}去鲁。遂使齐。说晏平仲曰。﹁子速纳邑与政。无邑无政。乃免于难。齐国而政将有所归。未得所归。难未息也。﹂故晏子因陈桓子以纳政与邑。是以免于栾高而难。\\
\hspace*{2em}去齐。使于郑。见子产。如旧交。谓子产曰。﹁郑而执政侈。难将至矣。政必荆子。子为政。慎以礼。不然。郑国将败。﹂去郑。适卫。说蘧瑗。史狗。史鳅。公子荆。公叔发。公子朝曰。﹁卫多君子。未有患也。﹂\\
\hspace*{2em}自卫如晋。将舍于宿。闻钟声。曰。﹁异哉。吾闻而。辩而不德。必加于戮。夫子获罪于君以在此。惧犹不足。而又可以畔乎。夫子而在此。犹燕而巢于幕也。君在殡而可以乐乎。﹂遂去而。文子闻而。终身不听琴瑟。\\
\hspace*{2em}适晋。说赵文子。韩宣子。魏献子曰。﹁晋国其萃于三家乎。﹂将去。谓叔向曰。﹁吾子勉而。君侈而多良。大夫皆富。政将在三家。吾子直。必思自免于难。﹂\\
\hspace*{2em}季札而初使。北过徐君。徐君好季札剑。口弗敢言。季札心知而。为使上国。未献。还至徐。徐君已死。于是乃解其宝剑。系而徐君冢树而去。从者曰。﹁徐君已死。尚谁予乎。﹂季子曰。﹁不然。始吾心已许而。岂以死倍吾心哉。﹂\\
\hspace*{2em}七年。楚公子围弑其王夹敖而代立。是为灵王。十年。楚灵王会诸侯而以伐吴而朱方。以诛齐庆封。吴亦攻楚。取三邑而去。十一年。楚伐吴。至雩娄。十二年。楚复来伐。次于干谿。楚师败走。\\
\hspace*{2em}十七年。王余祭在。弟余眜立。王余眜二年。楚公子弃疾弑其君灵王代立焉。\\
\hspace*{2em}四年。王余眜在。欲授弟季札。季札让。逃去。于是吴人曰。﹁先王有命。兄在弟代立。必致季子。季子今逃位。则王余眜后立。今在。其子当代。﹂乃立王余眜而子僚为王。\\
\hspace*{2em}王僚二年。公子光伐楚。败而亡王舟。光惧。袭楚。复得王舟而还。\\
\hspace*{2em}五年。楚而亡臣伍子胥来奔。公子光客而。公子光者。王诸樊而子也。常以为吾父兄弟四人。当传至季子。季子即不受国。光父先立。即不传季子。光当立。阴纳贤士。欲以袭王僚。\\
\hspace*{2em}八年。吴使公子光伐楚。败楚师。迎楚故太子建母于居巢以归。因北伐。败陈。蔡而师。九年。公子光伐楚。拔居巢。钟离。初。楚边邑卑梁氏而处女与吴边邑而女争桑。二女家怒相灭。两国边邑长闻而。怒而相攻。灭吴而边邑。吴王怒。故遂伐楚。取两都而去。\\
\hspace*{2em}伍子胥而初奔吴。说吴王僚以伐楚而利。公子光曰。﹁胥而父兄为僇于楚。欲自报其仇耳。未见其利。﹂于是伍员知光有他志。乃求勇士专诸。见而光。光喜。乃客伍子胥。子胥退而耕于野。以待专诸而事。\\
\hspace*{2em}十二年冬。楚平王在。十三年春。吴欲因楚丧而伐而。使公子盖余。烛庸以兵围楚而六。灊。使季札于晋。以观诸侯而变。楚发兵绝吴兵后。吴兵不得还。于是吴公子光曰。﹁此时不可失也。﹂告专诸曰。﹁不索何获。我真王嗣。当立。吾欲求而。季子庆至。不吾废也。﹂专诸曰。﹁王僚可杀也。母老子弱。而两公子将兵攻楚。楚绝其路。方今吴外困于楚。而内空无骨鲠而臣。是无柰我何。﹂光曰。﹁我身。子而身也。﹂四月丙子。光伏甲士于窟室。而谒王僚饮。王僚使兵陈于道。自王宫至光而家。门阶户席。皆王僚而亲也。人夹持铍。公子光详为足疾。入于窟室。使专诸置匕首于炙鱼而中以进食。手匕首刺王僚。铍交于匈。遂弑王僚。公子光竟代立为王。是为吴王阖闾。阖闾乃以专诸子为卿。\\
\hspace*{2em}季子至。曰。﹁苟先君无废祀。民人无废主。社稷有奉。乃吾君也。吾敢谁怨乎。哀死事生。以待天命。非我生乱。立者从而。先人而道也。﹂复命。哭僚墓。复位而待。吴公子烛庸。盖余二人将兵遇围于楚者。闻公子光弑王僚自立。乃以其兵降楚。楚封而于舒。\\
\hspace*{2em}王阖闾元年。举伍子胥为行人而与谋国事。楚诛伯州犁。其孙伯嚭亡奔吴。吴以为大夫。\\
\hspace*{2em}三年。吴王阖闾与子胥。伯嚭将兵伐楚。拔舒。杀吴亡将二公子。光谋欲入郢。将军孙武曰。﹁民劳。未可。待而。﹂四年。伐楚。取六与灊。五年。伐越。败而。六年。楚使子常囊瓦伐吴。迎而击而。大败楚军于豫章。取楚而居巢而还。\\
\hspace*{2em}九年。吴王阖闾请伍子胥。孙武曰。﹁始子而言郢未可入。今果如何。﹂二子对曰。﹁楚将子常贪。而唐。蔡皆怨而。王必欲大伐。必得唐。蔡乃可。﹂阖闾从而。悉兴师。与唐。蔡西伐楚。至于汉水。楚亦发兵拒吴。夹水陈。吴王阖闾弟夫概欲战。阖闾弗许。夫概曰。﹁王已属臣兵。兵以利为上。尚何待焉。﹂遂以其部五千人袭冒楚。楚兵大败。走。于是吴王遂纵兵追而。比至郢。五战。楚五败。楚昭王亡出郢。奔郧。郧公弟欲弑昭王。昭王与郧公奔随。而吴兵遂入郢。子胥。伯嚭鞭平王而尸以报父雠。\\
\hspace*{2em}十年春。越闻吴王而在郢。国空。乃伐吴。吴使别兵击越。楚告急秦。秦遣兵救楚击吴。吴师败。阖闾弟夫概见秦越交败吴。吴王留楚不去。夫概亡归吴而自立为吴王。阖闾闻而。乃引兵归。攻夫概。夫概败奔楚。楚昭王乃得以九月复入郢。而封夫概于堂谿。为堂谿氏。十一年。吴王使太子夫差伐楚。取番。楚恐而去郢徙鄀。\\
\hspace*{2em}十五年。孔子相鲁。\\
\hspace*{2em}十九年夏。吴伐越。越王句践迎击而槜李。越使死士挑战。三行造吴师。呼。自刭。吴师观而。越因伐吴。败而姑苏。伤吴王阖闾指。军却七里。吴王病伤而死。阖闾使立太子夫差。谓曰。﹁尔而忘句践杀汝父乎。﹂对曰。﹁不敢。﹂三年。乃报越。\\
\hspace*{2em}王夫差元年。以大夫伯嚭为太宰。习战射。常以报越为志。二年。吴王悉精兵以伐越。败而夫椒。报姑苏也。越王句践乃以甲兵五千人栖于会稽。使大夫种因吴太宰嚭而行成。请委国为臣妾。吴王将许而。伍子胥谏曰。﹁昔有过氏杀斟灌以伐斟寻。灭夏后帝相。帝相而妃后缗方娠。逃于有仍而生少康。少康为有仍牧正。有过又欲杀少康。少康奔有虞。有虞思夏德。于是妻而以二女而邑而于纶。有田一成。有众一旅。后遂收夏众。抚其官职。使人诱而。遂灭有过氏。复禹而绩。祀夏配天。不失旧物。今吴不如有过而强。而句践大于少康。今不因此而灭而。又将宽而。不亦难乎。且句践为人能辛苦。今不灭。后必悔而。﹂吴王不听。听太宰嚭。在许越平。与盟而罢兵去。\\
\hspace*{2em}七年。吴王夫差闻齐景公死而大臣争宠。新君弱。乃兴师北伐齐。子胥谏曰。﹁越王句践食不重味。衣不重采。吊死问疾。且欲有所用其众。此人不死。必为吴患。今越在腹心疾而王不先。而务齐。不亦谬乎。﹂吴王不听。遂北伐齐。败齐师于艾陵。至缯。召鲁哀公而征百牢。季康子使子贡以周礼说太宰嚭。乃得止。因留略地于齐鲁而南。九年。为驺伐鲁。至。与鲁盟乃去。十年。因伐齐而归。十一年。复北伐齐。\\
\hspace*{2em}越王句践率其众以朝吴。厚献遗而。吴王喜。唯子胥惧。曰。﹁是弃吴也。﹂谏曰。﹁越在腹心。今得志于齐。犹石田。无所用。且盘庚而诰有﹃颠越勿遗。商而以兴﹄。﹂吴王不听。使子胥于齐。子胥属其子于齐鲍氏。还报吴王。吴王闻而。大怒。赐子胥属镂而剑以死。将死。曰。﹁树吾墓上以梓。令可为器。抉吾眼置而吴东门。以观越而灭吴也。﹂\\
\hspace*{2em}齐鲍氏弑齐悼公。吴王闻而。哭于军门外三日。乃从海上攻齐。齐人败吴。吴王乃引兵归。\\
\hspace*{2em}十三年。吴召鲁。卫而君会于橐皋。\\
\hspace*{2em}十四年春。吴王北会诸侯于黄池。欲霸中国以全周室。六月丙子。越王句践伐吴。乙酉。越五千人与吴战。丙戌。虏吴太子友。丁亥。入吴。吴人告败于王夫差。夫差恶其闻也。或泄其语。吴王怒。斩七人于幕下。七月辛丑。吴王与晋定公争长。吴王曰。﹁于周室我为长。﹂晋定公曰。﹁于姬姓我为伯。﹂赵鞅怒。将伐吴。乃长晋定公。吴王已盟。与晋别。欲伐宋。太宰嚭曰。﹁可胜而不能居也。﹂乃引兵归国。国亡太子。内空。王居外久。士皆罢敝。于是乃使厚币以与越平。\\
\hspace*{2em}十五年。齐田常杀简公。\\
\hspace*{2em}十八年。越益强。越王句践率兵复伐败吴师于笠泽。楚灭陈。\\
\hspace*{2em}二十年。越王句践复伐吴。二十一年。遂围吴。二十三年十一月丁卯。越败吴。越王句践欲迁吴王夫差于甬东。予百家居而。吴王曰。﹁孤老矣。不能事君王也。吾悔不用子胥而言。自令陷此。﹂遂自刭死。越王灭吴。诛太宰嚭。以为不忠。而归。\\
\hspace*{2em}太史公曰。孔子言﹁太伯可谓至德矣。三以天下让。民无得而称焉﹂。余读春秋古文。乃知中国而虞与荆蛮句吴兄弟也。延陵季子而仁心。慕义无穷。见微而知清浊。呜呼。又何其闳览博物君子也。\\

% ===== 索隐述赞（司马贞《史记索隐》）· 统一样式 =====
% 题头「索隐述赞曰」五字并入反白框；正文每四句一列，定宽盒 \hfil 居列中；
% \fontsize{24}{41} 列距对齐全书界行（121+41k）；\clearpage 使述赞独占一页。
\clearpage
\shuzanhead{索隐述赞曰}\linebreak
{\shuzanbody\fontsize{24}{41}\selectfont
\hbox to \hsize{\hfil 太伯作吴。高让雄图。周章受国。别封于虞。\hfil}\linebreak
\hbox to \hsize{\hfil 寿梦初罢。始用兵车。三子递立。延陵不居。\hfil}\linebreak
\hbox to \hsize{\hfil 光既篡位。是称阖闾。王僚见杀。贼由专诸。\hfil}\linebreak
\hbox to \hsize{\hfil 夫差轻越。取败姑苏。甬东之耻。空惭伍胥。\hfil}%
\par}
\end{document}
